% GENERATED FILE. Edit canonical lessons, then run npm run generate:books.
% canonical-source-hash: fnv1a64:bd887cf902b63697
% canonical-lessons: SA-C133-nicah, SA-C133-unnatah, SA-C133-uddayate, SA-C133-tarati, SA-C133-kurdati

\chapter{Low, High, Flies, Swims, Jumps}
\label{ch:sa-low-high-flies-swims-jumps}
\hlchaptermodality{133}

\begin{chapteropening}
\textbf{By the end of this chapter:} \emph{I can say low, high, flies, swims and jumps.}

Complete the last lesson of chapter 133: I can say low, high, flies, swims and jumps.
\end{chapteropening}

\section[\texorpdfstring{\sk{नीचः} (nīcaḥ)}{नीचः (nīcaḥ)}]{\sk{नीचः} (nīcaḥ) — low}

\label{lesson:SA-C133-nicah}



\begin{quote}
Before the new one: say the Sanskrit for full, then the Sanskrit for empty.
\end{quote}

\subsection*{You'll want to know: \sk{नीचः}}
\textbf{\sk{नीचः}} — \emph{nīcaḥ} — \textquotedblleft{}low\textquotedblright{}.

\begin{grammarlens}[title={What you've built}]
One more word for this chapter.
\end{grammarlens}

\subsection*{Your turn}
\begin{itemize}
\raggedright
  \item \emph{Say it:} \emph{nīcaḥ}
  \item \emph{Say it:} \emph{nīcaḥ}, once more
  \item \emph{Say it:} say \emph{nīcaḥ}
  \item \emph{Recall:} say the Sanskrit for full, then the Sanskrit for empty, then say \emph{nīcaḥ} again
\end{itemize}

\begin{tcolorbox}[breakable,colback=teal!4,colframe=teal!35!black,title={Before you move on}]
What does \sk{नीचः} mean? (\textquotedblleft{}Low\textquotedblright{}.) Say it once more.
\end{tcolorbox}
\section[\texorpdfstring{\sk{उन्नतः} (unnataḥ)}{उन्नतः (unnataḥ)}]{\sk{उन्नतः} (unnataḥ) — high, tall}

\label{lesson:SA-C133-unnatah}



\begin{quote}
Before the new one: say the Sanskrit for empty, then the Sanskrit for low.
\end{quote}

\subsection*{You'll want to know: \sk{उन्नतः}}
\textbf{\sk{उन्नतः}} — \emph{unnataḥ} — \textquotedblleft{}high, tall\textquotedblright{}.

\begin{grammarlens}[title={What you've built}]
One more word for this chapter.
\end{grammarlens}

\subsection*{Your turn}
\begin{itemize}
\raggedright
  \item \emph{Say it:} \emph{unnataḥ}
  \item \emph{Say it:} \emph{unnataḥ}, once more
  \item \emph{Say it:} say \emph{unnataḥ}
  \item \emph{Recall:} say the Sanskrit for empty, then the Sanskrit for low, then say \emph{unnataḥ} again
\end{itemize}

\begin{tcolorbox}[breakable,colback=teal!4,colframe=teal!35!black,title={Before you move on}]
What does \sk{उन्नतः} mean? (\textquotedblleft{}High, tall\textquotedblright{}.) Say it once more.
\end{tcolorbox}
\section[\texorpdfstring{\sk{उड्डयते} (uḍḍayate)}{उड्डयते (uḍḍayate)}]{\sk{उड्डयते} (uḍḍayate) — flies}

\label{lesson:SA-C133-uddayate}



\begin{quote}
Before the new one: say the Sanskrit for low, then the Sanskrit for high.
\end{quote}

\subsection*{You'll want to know: \sk{उड्डयते}}
\textbf{\sk{उड्डयते}} — \emph{uḍḍayate} — \textquotedblleft{}flies\textquotedblright{}.

\begin{grammarlens}[title={What you've built}]
One more word for this chapter.
\end{grammarlens}

\subsection*{Your turn}
\begin{itemize}
\raggedright
  \item \emph{Say it:} \emph{uḍḍayate}
  \item \emph{Say it:} \emph{uḍḍayate}, once more
  \item \emph{Say it:} say \emph{uḍḍayate}
  \item \emph{Recall:} say the Sanskrit for low, then the Sanskrit for high, then say \emph{uḍḍayate} again
\end{itemize}

\begin{tcolorbox}[breakable,colback=teal!4,colframe=teal!35!black,title={Before you move on}]
What does \sk{उड्डयते} mean? (\textquotedblleft{}Flies\textquotedblright{}.) Say it once more.
\end{tcolorbox}
\section[\texorpdfstring{\sk{तरति} (tarati)}{तरति (tarati)}]{\sk{तरति} (tarati) — swims, crosses}

\label{lesson:SA-C133-tarati}



\begin{quote}
Before the new one: say the Sanskrit for high, then the Sanskrit for flies.
\end{quote}

\subsection*{You'll want to know: \sk{तरति}}
\textbf{\sk{तरति}} — \emph{tarati} — \textquotedblleft{}swims, crosses\textquotedblright{}.

\begin{grammarlens}[title={What you've built}]
One more word for this chapter.
\end{grammarlens}

\subsection*{Your turn}
\begin{itemize}
\raggedright
  \item \emph{Say it:} \emph{tarati}
  \item \emph{Say it:} \emph{tarati}, once more
  \item \emph{Say it:} say \emph{tarati}
  \item \emph{Recall:} say the Sanskrit for high, then the Sanskrit for flies, then say \emph{tarati} again
\end{itemize}

\begin{tcolorbox}[breakable,colback=teal!4,colframe=teal!35!black,title={Before you move on}]
What does \sk{तरति} mean? (\textquotedblleft{}Swims, crosses\textquotedblright{}.) Say it once more.
\end{tcolorbox}
\section[\texorpdfstring{\sk{कूर्दति} (kūrdati)}{कूर्दति (kūrdati)}]{\sk{कूर्दति} (kūrdati) — jumps}

\label{lesson:SA-C133-kurdati}



\begin{quote}
Before the new one: say the Sanskrit for flies, then the Sanskrit for swims.
\end{quote}

\subsection*{You'll want to know: \sk{कूर्दति}}
\textbf{\sk{कूर्दति}} — \emph{kūrdati} — \textquotedblleft{}jumps\textquotedblright{}.

\begin{grammarlens}[title={What you've built}]
One more word for this chapter.
\end{grammarlens}

\subsection*{Your turn}
\begin{itemize}
\raggedright
  \item \emph{Say it:} \emph{kūrdati}
  \item \emph{Say it:} \emph{kūrdati}, once more
  \item \emph{Say it:} say \emph{kūrdati}
  \item \emph{Recall:} say the Sanskrit for flies, then the Sanskrit for swims, then say \emph{kūrdati} again
\end{itemize}

\begin{tcolorbox}[breakable,colback=teal!4,colframe=teal!35!black,title={Before you move on}]
What does \sk{कूर्दति} mean? (\textquotedblleft{}Jumps\textquotedblright{}.) Say it once more.
\end{tcolorbox}
